% SPDX-License-Identifier: LPPL-1.3c
% Copyright (C) 2026 Christophe Jorssen
% Author and Current Maintainer: Christophe Jorssen <christophe.jorssen@gmail.com>
% LPPL maintenance status: maintained
% This file is part of luafemm. See LICENSE and LICENSES.md for its terms.
\section{Material library and constitutive laws}
\label{sec:materials}
The embedded catalogue contains all 282 records from the pinned FEMM 5
magnetic material library: air, permanent magnets, soft magnetic materials,
nonmagnetic metals and winding wires. A material is selected by an ASCII
identifier, such as |armco-pure-iron|, |copper|, |aluminum-6061-t6| or
|cast-alnico-5-lng37|. See Appendix~\ref{sec:catalogue} for every choice.
This is the magnetic solver's catalogue; dielectric and thermal material
databases belong to other FEMM solvers and are not imported.

The names, folders, every numeric field and every B--H point are preserved.
Duplicate names receive distinct identifiers; no record is overwritten. Installed documents need neither the FEMM
clone nor its executable.

\subsection{Choosing a catalogue edition}
The default is the bundled FEMM 5 collection. An edition selects data, not a
solver executable or an interchange file version. A FEM or ANS import uses its
embedded material values. No native executable or network lookup is needed.

\begin{key}{/femm/problem/material edition=\meta{choice} (initially femm5)}
Problem-scoped PGF |.is choice| selector: |femm5| or |femm4|. Set it in the
initial |femm/problem| options, before resolving materials or registering regions.
|femm5| uses the bundled MIT collection; |femm4| requires the explicit file below.
Material styles using |library| resolve within that problem's edition.
Styles and |.append style| work normally; ordinary TeX grouping restores defaults
for the next picture. An operating-point or sliding-position model owns its
captured materials and does not accept a different edition as an override.
\end{key}
\begin{key}{/femm/problem/material file=\meta{filename} (initially empty)}
Required only for |material edition=femm4|: the user's original |matlib.dat|
from FEMM revision \nolinkurl{7d9e8ed90772ddb939a3256f976e231bcb5cca97}.
Relative paths use the engine's working directory. The reader checks the exact
pinned byte fingerprint before accepting passive numeric data; a different file,
missing file, unsupported edition or a nonempty file with |femm5| is an error.
This input is not executed. See |data/femm5/README.md| for SHA-256 checksums.
The legacy collection is not distributed with this package.
\end{key}

For an existing calculation, obtain that exact file and set
|material edition=femm4| and
|material file=my-original-matlib.dat| in its initial
problem options. Continue to use the old material identifiers explicitly,
including component |core material| and |armature material| where needed.
The bundled component defaults now use |armco-pure-iron|. The old |pure-iron|
name is deliberately not an alias: curves, identifiers and default currents can
change between editions. Switching editions is a physical model change and
invalidates the solution cache. Recheck the field and circuit excitation after
migration. Explicit style overrides do not mutate either catalogue.

\begin{codeexample}[width=0pt]
\tikzset{femm/materials/catalogue-core/.style={library=armco-pure-iron}}
\pgfkeys{/femm/problem/catalogue example/.style={material edition=femm5},
  /femm/problem/catalogue example/.append style={mesh size=2}}
\begin{tikzpicture}[femm/problem={catalogue example,xmin=-8,xmax=8,
    ymin=-6,ymax=6},femm/boundary={side=all,y coefficient=.01}]
  \filldraw[fill=gray!15,femm/region={material=catalogue-core}]
    (-4,-2) rectangle (4,2);
  \pic[femm/lines=9,femm/arrows=false] {femm field};
\end{tikzpicture}
\end{codeexample}

\subsection{Custom and derived materials}
Define a material as a style under |/tikz/femm/materials|. The style name
then becomes a valid |material| choice. Define styles before the first
use, and use a name that is not reserved by the catalogue.
\begin{codeexample}[code only]
\tikzset{
  femm/materials/my-iron/.style={
    bh={0/0,.5/80,1/180,1.3/350,1.5/900,1.8/7000,2/20000}},
  femm/materials/my-metal/.style={library=copper,mu r=1.01},
  femm/materials/recoil-magnet/.style={mu r=1.05,remanence=1.2},
  core/.style={femm/region={material=my-iron},draw,fill=gray!20}
}
\end{codeexample}

The following keys are evaluated while resolving a material style. Empty
internal slots mean ``inherit from the library, if any''. The effective
defaults below apply when there is no library base. Explicit properties
override the base, independently of their order in the style.

\begin{key}{/tikz/femm/materials/library=\meta{choice}}
Use a catalogue record as the base of a custom material. This is a PGF
|.is choice| handler, with no initial base. A private copy is created;
overrides never change the catalogue itself. A material name is model-wide:
do not redefine its physical law between regions of the same model.
\end{key}
\begin{key}{/tikz/femm/materials/mu r=\meta{positive number}}
Relative permeability of a linear isotropic material; effective default 1.
A nonempty |bh| table takes precedence. Setting |mu r| on a derived
nonlinear material does not remove its B--H table.
\end{key}
\begin{key}{/tikz/femm/materials/bh=\marg{B/H pairs}}
Comma-separated numeric pairs: B in tesla, H in A/m. The table must have
at least two points, begin with |0/0|, and have strictly increasing B and H.
There is no curve by default. A nonempty table replaces the base table;
an empty value inherits it. The package interpolates H(B) linearly. Beyond
the final point, the slope is $1/\mu_0$. This is deliberately not FEMM's
spline interpolation, so imported data do not imply bitwise FEMM results.
\end{key}
\begin{key}{/tikz/femm/materials/coercivity=\meta{nonnegative number}}
Prescribed source strength $H_c$ in A/m; effective default zero.
It can accompany either |mu r| or a \emph{shifted} B--H table. Its direction
is chosen by each region's |magnetization angle|.
\end{key}
\begin{key}{/tikz/femm/materials/remanence=\meta{nonnegative number}}
Remanent induction $B_r$ in tesla, for a linear recoil law; effective
default zero. A positive value is incompatible with a B--H table or
nonzero |coercivity|, including inherited ones. It gives
$H_c=B_r/(\mu_0\mu_r)$. This extrapolated intercept is not necessarily the
measured coercivity of a nonlinear hard magnet.
\end{key}

\subsection{Permanent magnets}
The constitutive convention is
\[
 \mathbf H=\nu(|\mathbf B|)\mathbf B-H_c\mathbf m,
 \qquad \mathbf m=(\cos\alpha,\sin\alpha),\qquad
 \nu(B)=H_{\rm table}(B)/B.
\]
Along the positive magnetisation direction, the physical field is
$H_{\rm physical}(B)=H_{\rm table}(B)-H_c$. Thus the table starts at
|0/0| while the physical demagnetisation branch starts at $H=-H_c$.
Remanence occurs where $H_{\rm table}(B_r)=H_c$. Do not paste a measured
negative-H demagnetisation table directly into |bh|.

The |cast-alnico-5-lng37| record has $H_c=48600$ A/m and a table endpoint
at $B=1.204$ T, $H_{\rm table}=48600$ A/m. The example
|examples/alnico-u.tex| uses that table with illustrative soft iron;
|examples/materials-profiles.tex| uses the catalogue's |armco-pure-iron| instead
(complete source and result in Section~\ref{ex:materials-profiles}).
Both produce a field without imposed current. The magnetisation direction
is prescribed, not an internal history variable: hysteresis, irreversible
demagnetisation and measured transverse anisotropy are not modelled.

\subsection{Example: explicit material tables}
The following U magnet shows the original illustrative soft-iron curve
and the shifted alnico table explicitly in reusable material styles.
The local air-gap refinement regions are declared after the solids and
occupy only the gaps. The output includes field lines, geometry and the
computed diagnostics. These illustrative iron data are not a commercial
steel specification.
\luafemmexamplebegin{alnico-u}

% manual example: alnico-u/1

\begin{codeexample}[setup code]
\definecolor{alnico}{HTML}{B75046}
\definecolor{flux}{HTML}{147B83}
\tikzset{
  femm/materials/iron/.style={bh={0/0,.5/80,1/180,1.3/350,
    1.5/900,1.65/2500,1.8/7000,2/20000}},

  femm/materials/alnico/.style={coercivity=48600,
    bh={0/0,.2781/1600,.6383/4100,.8234/6300,.9214/8600,
      1.013/13400,1.064/18200,1.106/24300,1.149/33000,
      1.184/42530,1.204/48600}},
  iron/.style={femm/region={material=iron},draw=black!65,fill=black!9},
  magnet/.style={femm/region={material=alnico,magnetization angle=180},
    draw=alnico!80!black,fill=alnico!25}
}
\end{codeexample}

Apply these styles to the physical paths before requesting the field.

% manual example: alnico-u/2

\begin{codeexample}[width=0pt]
\begin{tikzpicture}[femm/problem={xmin=-100,xmax=100,ymin=-85,ymax=100,
  mesh size=2,mesher=delaunay}]
  \filldraw[iron] (-40,-30)--(40,-30)--(40,30)--(25,30)
    --(25,-15)--(-25,-15)--(-25,30)--(-40,30)--cycle;
  \filldraw[magnet] (-25,-30) rectangle (25,-15);
  \filldraw[iron] (-40,31) rectangle (40,46);

  \path[femm/region={mesh size=.5}] (-40,30) rectangle (-25,31);
  \path[femm/region={mesh size=.5}] (25,30) rectangle (40,31);
  \begin{scope}
    \clip (-61,-45) rectangle (62,64);
    \pic[femm/lines=29,femm/arrow spacing=26mm,
      femm/field style={flux}] {femm field};
  \end{scope}
  \node[fill=white,inner sep=2pt,font=\small] at (0,39) {Soft-iron armature};
  \node[fill=white,inner sep=2pt,font=\small] at (0,4) {Air: $J_z=0$ everywhere};
  \node[fill=white,inner sep=2pt,text=flux] at (-32.5,27) {N};
  \node[fill=white,inner sep=2pt,text=flux] at (32.5,27) {S};
  \node[fill=alnico!25,inner sep=2pt,font=\small] at (0,-20) {Alnico 5 (LNG37)};
  \draw[white,line width=2.5pt] (15,-25)--(-15,-25);
  \draw[-Stealth,line width=1pt,alnico!80!black] (15,-25)--(-15,-25)
    node[midway,below,fill=alnico!25,inner sep=1pt,font=\scriptsize] {magnetization direction};
  \draw[thin] (40,30)--(50,30) (40,31)--(50,31);
  \draw[thin,-Stealth] (48,25)--(48,30);
  \draw[thin,-Stealth] (48,36)--(48,31);
  \node[anchor=west,fill=white,font=\small,inner sep=1pt] at (50,30.5) {$e=1\,\mathrm{mm}$};
  \draw[thin] (-25,-30)--(-25,-38) (25,-30)--(25,-38);
  \draw[Stealth-Stealth,thin] (-25,-36)--(25,-36)
    node[midway,fill=white,font=\small,inner sep=2pt] {$50\,\mathrm{mm}$};
  \node[fill=white,font=\small,text=flux,inner sep=2pt]
    at (0,58) {$\mathbf B$ contours, oriented from N to S in air};
\end{tikzpicture}
\end{codeexample}

\luafemmexampleend{alnico-u}

\subsection{Static interpretation of the catalogue}
All records in this revision have equal |Mu_x| and |Mu_y|. For |LamType=0|,
a fill fraction $f$ is applied as
\[
 \mu_{r,\rm eff}=f\mu_r+(1-f),\qquad
 B_{\rm eff}(H)=f B(H)+(1-f)\mu_0H.
\]
This accounts for the ten records with fill fraction 0.98. Winding records have unit relative permeability. Type 3 round-wire records
select the explicit AC homogenisation law and require positive frequency;
for a DC winding use an ordinary conductor material and specify its copper
fill. Type 8 copper-clad aluminium retains only its static scalar interpretation. Custom linear anisotropy and directional sheet mixtures are described in
Section~\ref{sec:static}. The catalogue adapter retains its audited isotropic
interpretation and rejects unsupported database combinations.

Conductivity is used by the named DC winding model in
Section~\ref{sec:static} and by the harmonic diffusion solver in
Section~\ref{sec:harmonic}. Section~\ref{sec:losses} adds supported linear principal loss angles and
insulated-sheet thickness/fill. Nonlinear catalogue B/H curves retain their
static meaning; selecting them does not create a nonlinear harmonic law.
Supported round-wire diameter and strand data select the effective laws of
Section~\ref{sec:harmonic-material-models}, including square and CCA laws. The real material source |J_re| is converted from A/mm$^2$ to A/m$^2$.
The bundled FEMM 5 copper-clad aluminium records have zero imposed current.
The optional legacy FEMM 4 collection contains 69 wire records with a default
of 1 A/mm$^2$; explicitly use |current density=0| to make those regions passive.
CCA records are available as nonmagnetic conductors in DC, but their special
harmonic wire laws are not implemented. Stored fields do not imply support for
all native FEMM models.

\subsection{Provenance and data terms}
The bundled source is |femm-qt/matlibs/matlib.dat|, revision
\nolinkurl{31ee4bf6a8cae1a1860c108e7934ab3403bfd6e3}.
The unchanged source, MIT notice and material-library explanation are in
|data/femm5/|. The generated Lua transcription and this appendix retain those
MIT terms; the original package code and prose remain LPPL 1.3c or later.
The solver's separate GPL licence does not apply to this expressly MIT-licensed
collection. No solver code, binary or provenance paper is bundled.
FEMM 4's AFPL material collection is excluded from the public distribution.
